% ===== references =====
\begin{thebibliography}{99}
\bibitem{wakili2025} A. A. Wakili, S. Guni, S. A. Khan, W. Yu, and W. Jung, ``Advancing machine learning strategies for power consumption-based IoT botnet detection,'' \textit{Sensors}, vol. 25, p. 7553, 2025.
\bibitem{sudheera2026} K. L. K. Sudheera \textit{et al.}, ``Federated learning and data mining-based botnet attack detection framework for Internet of Things,'' \textit{Sensors}, vol. 26, p. 1573, 2026.
\bibitem{memos2026} V. A. Memos, C. L. Stergiou, A. I. Bermperis, A. P. Plageras, and K. E. Psannis, ``A novel architecture for mitigating botnet threats in AI-powered IoT environments,'' \textit{Sensors}, vol. 26, p. 572, 2026.
\bibitem{chai2026} J. Chai, Y. Shen, W. Tian, X. Sun, and L. Tian, ``RI graph sampling: A hybrid graph sampling method for large-scale botnet detection,'' \textit{Sensors}, vol. 26, p. 6154, 2026.
\bibitem{georgiades2025} M. Georgiades and F. Hussain, ``An explainable AI approach for interpretable cross-layer intrusion detection in Internet of Medical Things,'' \textit{Electronics}, vol. 14, p. 3218, 2025.
\bibitem{abujazoh2025} M. Abu-Jazoh, I. Almomani, and K. E. Sabri, ``DCmal-2025: A novel routing-based disconnectivity malware---development, impact, and countermeasures,'' \textit{Appl. Sci.}, vol. 15, p. 10219, 2025.
\bibitem{banoth2026} R. Banoth \textit{et al.}, ``Evolving IoT botnet threats and practical honeypot observation: A summary review and experimental study,'' \textit{J. Cybersecur. Priv.}, vol. 6, p. 82, 2026.
\bibitem{babenko2026} T. Babenko \textit{et al.}, ``Hybrid GNN--LSTM architecture for probabilistic IoT botnet detection with calibrated risk assessment,'' \textit{Computers}, vol. 15, p. 26, 2026.
\bibitem{aljaafreh2026} M. Alja'afreh, A. Karime, and A. Oukaira, ``Privacy-aware and resource-efficient split learning for IoT botnet detection: A multi-dataset experimental and systems evaluation,'' \textit{IoT}, vol. 7, p. 78, 2026.
\bibitem{awad2026} L. Awad, S. Al-Haj Baddar, and A. Sleit, ``Robust and noise-resilient botnet detection framework using heterogeneous radial basis function neural network,'' \textit{Appl. Sci.}, vol. 16, p. 3379, 2026.
\bibitem{hua2026} M. Hua, X. Zhang, and Y. Li, ``SHAP-guided feature selection for IoT botnet detection: A comparative study with conventional methods,'' \textit{Appl. Sci.}, vol. 16, p. 8840, 2026.
\bibitem{antonakakis2017} M. Antonakakis \textit{et al.}, ``Understanding the Mirai botnet,'' in \textit{Proc. 26th USENIX Security Symp.}, 2017, pp. 1093--1110.
\bibitem{iot23} S. Garcia, A. Parmisano, and M. J. Erquiaga, ``IoT-23: A labeled dataset with malicious and benign IoT network traffic,'' Zenodo, 2020, doi: 10.5281/zenodo.4743746.
\bibitem{chawla2002} N. V. Chawla, K. W. Bowyer, L. O. Hall, and W. P. Kegelmeyer, ``SMOTE: Synthetic minority over-sampling technique,'' \textit{J. Artif. Intell. Res.}, vol. 16, pp. 321--357, 2002.
\bibitem{breiman2001} L. Breiman, ``Random forests,'' \textit{Mach. Learn.}, vol. 45, no. 1, pp. 5--32, 2001.
\bibitem{chen2016} T. Chen and C. Guestrin, ``XGBoost: A scalable tree boosting system,'' in \textit{Proc. 22nd ACM SIGKDD}, 2016, pp. 785--794.
\bibitem{ke2017} G. Ke \textit{et al.}, ``LightGBM: A highly efficient gradient boosting decision tree,'' in \textit{Proc. NeurIPS}, vol. 30, 2017.
\bibitem{kingma2015} D. P. Kingma and J. Ba, ``Adam: A method for stochastic optimization,'' in \textit{Proc. ICLR}, 2015.
\bibitem{lundberg2017} S. M. Lundberg and S.-I. Lee, ``A unified approach to interpreting model predictions,'' in \textit{Proc. NeurIPS}, vol. 30, 2017, pp. 4765--4774.
\bibitem{lundberg2020} S. M. Lundberg \textit{et al.}, ``From local explanations to global understanding with explainable AI for trees,'' \textit{Nat. Mach. Intell.}, vol. 2, pp. 56--67, 2020.
\bibitem{pedregosa2011} F. Pedregosa \textit{et al.}, ``Scikit-learn: Machine learning in Python,'' \textit{J. Mach. Learn. Res.}, vol. 12, pp. 2825--2830, 2011.
\bibitem{lemaitre2017} G. Lema\^{i}tre, F. Nogueira, and C. K. Aridas, ``Imbalanced-learn: A Python toolbox to tackle the curse of imbalanced datasets in machine learning,'' \textit{J. Mach. Learn. Res.}, vol. 18, no. 17, pp. 1--5, 2017.
\bibitem{paszke2019} A. Paszke \textit{et al.}, ``PyTorch: An imperative style, high-performance deep learning library,'' in \textit{Proc. NeurIPS}, vol. 32, 2019.
\end{thebibliography}